\section{Fixed middle data for selected windows}

Let the chain have sites $\{0,\ldots,L-1\}$ and local space $\C^q$.
All intervals below are intersected with this chain. For nonnegative
integers $a,b,\ell$ and positions $0\le p,s<2\ell$, put
\begin{align}
    K & =[a+2\ell,b+1),                      & Z &
    =[a+2\ell+1,b),\notag                                         \\
    T & =[0,a+p+1),                          & E &
    =[a+p+1,b+s+1),\notag                                         \\
    R & =\{0,\ldots,L-1\}\setminus(T\cup E), & C & =E\setminus Z.
    \label{eq:selected_middle_regions}
\end{align}
These statements provide supporting locality algebra for
\cite[Proposition~5.7]{Arad2016RigorousRG}.

\begin{lemma}[The central regions lie in the selected middle]
    \label{thm:selected_middle_core_inclusions}
    \lean{TNLean.MPS.AreaLaw.core_and_strict_core_subset_selected_middle}
    \leanok
    For the data above, $K\subseteq E$ and $Z\subseteq E$.
    No separation assumption on $a,b$ is needed for these inclusions.
\end{lemma}
\begin{proof}\leanok
    The inequalities $p<2\ell$ and $s\ge0$ show that the left endpoint of
    $E$ is no greater than that of $K$, while its right endpoint is no
    smaller. The same comparison applies to $Z$. Monotonicity of clipped
    intervals proves the assertions.
\end{proof}

\begin{lemma}[Cardinality of the selected collars]
    \label{thm:selected_middle_collar_cardinality}
    \lean{TNLean.MPS.AreaLaw.card_selected_middle_sdiff_strict_core_le}
    \leanok
    If $a+2\ell\le b$, then $|C|\le4\ell$.
\end{lemma}
\begin{proof}\leanok
    The complement of the strict core within $E$ is
    \begin{align}
        C=[a+p+1,a+2\ell+1)\cup
        [b,b+s+1).
        \label{eq:selected_middle_collar_union}
    \end{align}
    The two intervals may overlap when the windows touch. Their cardinalities
    are at most $2\ell-p$ and $s+1$, respectively. Both are at most $2\ell$,
    so the cardinality of their union is at most $4\ell$.
\end{proof}

\begin{lemma}[Middle and crossing factors chosen before the exterior operators]
    \label{thm:selected_middle_fixed_physical_data}
    \lean{TNLean.MPS.AreaLaw.exists_physical_selected_window_fixed_middle_data}
    \leanok
    \uses{thm:selected_middle_core_inclusions,thm:three_region_partial_support,
        thm:three_region_middle_lift_equality,thm:three_region_middle_embedding,
        thm:three_region_two_window_collars,thm:three_region_weighted_middle_remainder,
        thm:two_marker_prescribed_span}
    Assume $a+2\ell\le b$ and $b+2\ell\le L-1$.
    Let $W^L_r,W^R_r$, for $0\le r<2\ell$, be operators on the full chain
    supported on $[a+r,a+r+2)$ and
    $[b+r,b+r+2)$, respectively.
    Let $M$ be an arbitrary operator on $K$. Write $P_M$ for its canonical
    extension to the full chain and $M_E$ for its canonical extension to $E$.

    Use labels $L_r,R_r,B_0,B_1,B_2$, and let $O$ be all labels except
    $L_p,R_s$. There exist sites $e_L,e_R\in E$, a middle family $H^M$,
    crossing operators $U,V_0$, and exterior families
    $A_z\in\operatorname{End}(\mathcal H_T)$,
    $B_z\in\operatorname{End}(\mathcal H_R)$, indexed by
    $z\in\{0,\ldots,q^2-1\}$, with the following properties.

    First, the sites are the literal selected endpoints:
    \begin{align}
        e_L=a+p+1,\qquad e_R=b+s,
        \label{eq:selected_middle_endpoints}
    \end{align}
    and the central middle factor is $H^M_{B_1}=M_E$.
    For $z$ corresponding to a pair $(i,j)\in\{0,\ldots,q-1\}^2$,
    let $N_e(z)$ be the one-site matrix unit $|i\rangle\langle j|$
    canonically extended from $e$ to $E$.
    In three-region coordinates the actual selected crossings satisfy
    \begin{align}
        (W^L_p)^g & =U\otimes I_R,                 &
        (W^R_s)^g & =I_T\otimes V_0,\notag           \\
        U         & =\sum_z A_z\otimes N_{e_L}(z), &
        V_0       & =\sum_z N_{e_R}(z)\otimes B_z.
        \label{eq:selected_middle_fixed_crossings}
    \end{align}
    Every $N_{e_L}(z)$ and $N_{e_R}(z)$ is supported on $C$.
    Moreover, for every complex weight family $z_i$ there is an operator
    $K_z$ supported on $C$ such that
    \begin{align}
        \sum_{i\in O}z_iH^M_i=z_{B_1}M_E+K_z.
        \label{eq:selected_middle_weighted_sum}
    \end{align}

    All the preceding choices are made before the following exterior data.
    For every pair of full-chain operators $P_L,P_R$ supported on
    $[0,a+1)$ and
    $[b+2\ell,L)$, respectively, set
    \begin{align}
        F_{L_r}=W^L_r,\quad F_{R_r}=W^R_r,\quad
        F_{B_0}=P_L,\quad F_{B_1}=P_M,\quad F_{B_2}=P_R.
        \label{eq:selected_middle_full_family}
    \end{align}
    There are exterior regional families $H^L,H^R$ such that, for $i\in O$,
    \begin{align}
        F_i^g=H_i^L\otimes I_E\otimes I_R
        +I_T\otimes H_i^M\otimes I_R
        +I_T\otimes I_E\otimes H_i^R,
        \label{eq:selected_middle_late_exterior}
    \end{align}
    with $H^L_{B_1}=H^R_{B_1}=0$.
    The sites $e_L,e_R$ may coincide. No positive-local-dimension assumption
    is imposed.
\end{lemma}
\begin{proof}\leanok
    \uses{thm:selected_middle_core_inclusions,thm:three_region_partial_support,
        thm:three_region_middle_lift_equality,thm:three_region_middle_embedding,
        thm:three_region_two_window_collars,thm:three_region_weighted_middle_remainder,
        thm:two_marker_prescribed_span}
    Apply the three-region collar decomposition with both exterior bulk
    operators set to zero and with central operator $P_M$.
    The two selected bonds meet $E$ only at $a+p+1$ and $b+s$.
    Both endpoints lie in $C$, as follows directly from $p,s<2\ell$.

    Canonical placement of $M$ factors through its inclusion into $E$.
    Since the central exterior components vanish, the chosen central middle
    factor satisfies
    \begin{align}
        P_M^g=I_T\otimes H^M_{B_1}\otimes I_R
        =I_T\otimes M_E\otimes I_R.
        \label{eq:selected_middle_central_lift}
    \end{align}
    The endpoint $e_L$ supplies an actual site of $E$, so the two middle
    factors in (\ref{eq:selected_middle_central_lift}) coincide,
    including at $q=0$.

    The partial matrix elements of each crossing are supported on its single
    middle endpoint. An operator supported on one site belongs to the span of
    the canonically extended matrix units at that site. Retaining these
    prescribed spanning families in the two product decompositions gives
    (\ref{eq:selected_middle_fixed_crossings}). Their support lies in $C$.
    The weighted supported-remainder identity then gives
    (\ref{eq:selected_middle_weighted_sum}).

    Finally, insert arbitrary $P_L,P_R$ with the stated supports by changing
    only the corresponding exterior bulk components. The middle family and
    the exterior components of the window labels remain fixed. This proves
    (\ref{eq:selected_middle_late_exterior}) with the central exterior
    components equal to zero.
\end{proof}
